\documentclass[11pt,a4paper]{article}

\usepackage[utf8]{inputenc}
\usepackage[T1]{fontenc}
\usepackage[french]{babel}
\usepackage{geometry}
\usepackage{xcolor}
\usepackage{hyperref}
\usepackage{tabularx}
\usepackage{longtable}
\usepackage{enumitem}
\usepackage{listings}
\usepackage{titlesec}

\geometry{margin=2.4cm}
\hypersetup{
  colorlinks=true,
  linkcolor=blue!55!black,
  urlcolor=blue!55!black
}

\definecolor{chorusblue}{HTML}{1D4ED8}
\definecolor{chorusgray}{HTML}{F3F4F6}
\definecolor{chorusdark}{HTML}{111827}

\titleformat{\section}
  {\Large\bfseries\color{chorusblue}}
  {\thesection}{0.75em}{}

\titleformat{\subsection}
  {\large\bfseries\color{chorusdark}}
  {\thesubsection}{0.75em}{}

\lstdefinestyle{chorus}{
  backgroundcolor=\color{chorusgray},
  basicstyle=\ttfamily\small,
  breaklines=true,
  frame=single,
  rulecolor=\color{black!15},
  columns=fullflexible
}
\lstset{style=chorus}

\setlist[itemize]{leftmargin=1.2cm}
\setlist[enumerate]{leftmargin=1.2cm}

\title{
  \vspace{-1.5cm}
  \textbf{\Huge Chorus Pilot}\\
  \vspace{0.3cm}
  \Large Présentation de l'application et de ses fonctionnalités
}
\author{Projet Chorus}
\date{\today}

\begin{document}

\maketitle
\thispagestyle{empty}

\begin{center}
  \vspace{0.5cm}
  \begin{tabularx}{0.9\textwidth}{|X|}
    \hline
    \textbf{Résumé} \\
    Chorus Pilot est une application web destinée à préparer, contrôler,
    synchroniser et soumettre des factures vers Chorus Pro via PISTE. Elle combine
    une interface React, une API Python FastAPI, une base PostgreSQL et une stack
    Docker pour couvrir le cycle de vie opérationnel d'une facture fournisseur. \\
    \hline
  \end{tabularx}
\end{center}

\newpage
\tableofcontents
\newpage

\section{Contexte et objectifs}

Le projet \textbf{Chorus Pilot} répond au besoin d'industrialiser le
traitement des factures à destination de Chorus Pro. L'application facilite la
préparation des données, la vérification des référentiels, l'extraction de
contenu depuis les documents PDF et la soumission des factures via les API
Chorus.

Les objectifs principaux sont les suivants :

\begin{itemize}
  \item centraliser la consultation des factures synchronisées depuis Chorus ;
  \item préparer une facture à partir d'un PDF ou de données saisies ;
  \item extraire automatiquement les informations utiles d'un document ;
  \item contrôler les données avant soumission ;
  \item construire un payload compatible Chorus Pro ;
  \item transmettre la facture au backend, qui réalise l'upload PDF puis la
        soumission vers Chorus ;
  \item suivre les statuts, les retours API et l'activité technique.
\end{itemize}

\section{Vue d'ensemble fonctionnelle}

L'application est organisée autour de plusieurs écrans métier accessibles après
authentification.

\begin{longtable}{|p{0.25\textwidth}|p{0.68\textwidth}|}
  \hline
  \textbf{Module} & \textbf{Rôle} \\
  \hline
  Tableau de bord &
  Donne une vision synthétique des factures, montants, statuts et synchronisations.
  Il permet également de déclencher une synchronisation des fournisseurs et des
  factures. \\
  \hline
  Factures &
  Liste les factures connues de l'application avec recherche, filtres Chorus,
  pagination, tri, montants, statuts et accès au détail. \\
  \hline
  Nouvelle facture &
  Permet d'importer un PDF ou un XML, d'extraire les données, de compléter les
  informations nécessaires et de soumettre la facture à Chorus Pro. \\
  \hline
  Détail facture &
  Affiche les informations d'une facture, son historique, ses statuts, ses pièces
  jointes et, selon la configuration, les traces techniques utiles au diagnostic. \\
  \hline
  Référentiels &
  Donne accès aux clients publics et aux services destinataires. Le module permet
  la recherche, la consultation d'informations SIRET et la synchronisation de
  l'annuaire Chorus. \\
  \hline
  Paramètres &
  Centralise la configuration Chorus : environnement, connecteur, endpoint backend,
  contexte technique, modes de paiement, type de facture, TVA et test de connexion. \\
  \hline
  Connexion / Profil &
  Gère l'accès à l'application via une session JWT portée par le backend. \\
  \hline
\end{longtable}

\section{Workflow de soumission d'une facture}

Le flux nominal de traitement d'une facture se déroule en plusieurs étapes :

\begin{enumerate}
  \item l'utilisateur se connecte à l'application ;
  \item il ouvre le module \textit{Nouvelle facture} ;
  \item il importe un document PDF ou XML ;
  \item l'application extrait le texte et tente d'identifier les données utiles :
        numéro de facture, date, montants, TVA, fournisseur, destinataire et service ;
  \item l'utilisateur vérifie ou complète les champs obligatoires ;
  \item le frontend construit un payload métier normalisé ;
  \item le backend reçoit la demande de soumission ;
  \item le backend charge le PDF vers Chorus Pro ;
  \item le backend soumet la facture en référençant la pièce jointe principale ;
  \item le retour Chorus est affiché dans l'interface et conservé pour le suivi.
\end{enumerate}

\subsection{Modes de dépôt}

Le projet prend en compte les principaux modes de dépôt attendus par Chorus :

\begin{itemize}
  \item \textbf{DEPOT\_PDF\_API} : dépôt d'une facture PDF via API, avec upload
        préalable du fichier ;
  \item \textbf{SAISIE\_API} : soumission d'une facture sous forme de données
        structurées sans pièce jointe principale obligatoire ;
  \item \textbf{DEPOT\_PDF\_SIGNE\_API} : dépôt PDF signé, selon la configuration
        Chorus disponible.
\end{itemize}

\section{Architecture applicative}

Chorus Pilot repose sur une architecture web en plusieurs couches :

\begin{lstlisting}
Navigateur
  -> nginx frontend
    -> Application React statique
    -> Proxy /api/* vers FastAPI
      -> Authentification JWT
      -> Extraction PDF
      -> Client Chorus / PISTE
      -> API de données /api/db/*
        -> PostgreSQL
\end{lstlisting}

\subsection{Frontend}

Le frontend est développé avec React, TypeScript et Vite. Il utilise notamment :

\begin{itemize}
  \item \textbf{React Router} pour la navigation ;
  \item \textbf{TanStack Query} pour le chargement, le cache et la synchronisation
        des données ;
  \item \textbf{Tailwind CSS} et des composants UI pour l'interface ;
  \item \textbf{Recharts} pour les graphiques du tableau de bord ;
  \item des modules métier dans \texttt{src/lib/chorus/} pour préparer, valider
        et hydrater les payloads Chorus.
\end{itemize}

\subsection{Backend}

Le backend est une API Python FastAPI exposée par
\texttt{scripts/chorus\_submit\_server.py}. Il joue le rôle de passerelle entre le
navigateur, la base de données et Chorus Pro.

Ses responsabilités principales sont :

\begin{itemize}
  \item authentifier les utilisateurs ;
  \item protéger les secrets Chorus côté serveur ;
  \item récupérer les tokens OAuth2 ;
  \item extraire le texte des PDF ;
  \item synchroniser l'annuaire Chorus ;
  \item rechercher des factures fournisseur côté Chorus ;
  \item uploader les pièces jointes ;
  \item soumettre les factures à Chorus Pro ;
  \item exposer les données applicatives via \texttt{/api/db/*}.
\end{itemize}

\subsection{Base de données}

La base PostgreSQL conserve les référentiels, factures, historiques,
synchronisations et données applicatives. Les migrations sont gérées avec
Alembic. En environnement Docker, les données sont stockées dans un volume
persistant afin d'éviter toute perte lors des redéploiements.

\section{Principaux endpoints}

\begin{longtable}{|p{0.35\textwidth}|p{0.58\textwidth}|}
  \hline
  \textbf{Endpoint} & \textbf{Description} \\
  \hline
  \texttt{GET /health} &
  Vérifie la disponibilité de l'API. \\
  \hline
  \texttt{POST /api/auth/login} &
  Authentifie l'utilisateur et initialise la session. \\
  \hline
  \texttt{GET /api/auth/me} &
  Retourne les informations de la session courante. \\
  \hline
  \texttt{POST /api/chorus/oauth-token} &
  Récupère un token OAuth2 Chorus côté serveur. \\
  \hline
  \texttt{POST /api/extract-pdf-text} &
  Extrait le texte d'un PDF encodé en base64. \\
  \hline
  \texttt{POST /api/chorus/submit} &
  Réalise l'upload PDF puis la soumission de la facture vers Chorus. \\
  \hline
  \texttt{POST /api/chorus/annuaire/sync} &
  Synchronise l'annuaire Chorus. \\
  \hline
  \texttt{GET /api/db/*} &
  Expose les données applicatives stockées en base. \\
  \hline
\end{longtable}

\section{Configuration}

La configuration est séparée entre les variables serveur et les variables
visibles côté navigateur.

\subsection{Variables serveur}

Les secrets doivent rester côté serveur, dans \texttt{.env.local}. Les variables
importantes sont notamment :

\begin{lstlisting}
CHORUS_ENV=prod
CHORUS_PROD_CLIENT_ID=
CHORUS_PROD_CLIENT_SECRET=
CHORUS_PROD_TECHNICAL_LOGIN=
CHORUS_PROD_TECHNICAL_PASSWORD=
CHORUS_PROD_SUPPLIER_ID=
CHORUS_PROD_DESTINATAIRE_CODE=
CHORUS_PROD_SERVICE_EXECUTANT_CODE=
CHORUS_PROD_BANK_ACCOUNT_CODE=
\end{lstlisting}

Pour la qualification, l'application utilise les variantes
\texttt{CHORUS\_QUALIF\_*}.

\subsection{Variables frontend}

Les variables \texttt{VITE\_*} sont injectées dans le bundle navigateur et peuvent
être visibles par l'utilisateur. Elles ne doivent donc pas contenir de secrets de
production.

\section{Sécurité}

L'application applique plusieurs principes de sécurité :

\begin{itemize}
  \item les secrets Chorus ne sont pas exposés directement au navigateur ;
  \item le frontend passe par le backend pour les appels sensibles ;
  \item l'accès utilisateur repose sur une authentification avec session JWT ;
  \item le backend peut limiter les fournisseurs synchronisés via
        \texttt{CHORUS\_ALLOWED\_SUPPLIER\_IDS} ;
  \item les traces techniques sont utiles au diagnostic mais doivent être
        contrôlées en production.
\end{itemize}

Dans l'état actuel, l'application utilise un compte applicatif partagé. Une
évolution multi-utilisateurs complète nécessiterait une table dédiée aux
utilisateurs, des rôles et une traçabilité par personne.

\section{Synchronisation et référentiels}

La synchronisation permet de maintenir les données locales cohérentes avec
Chorus Pro. Elle concerne notamment :

\begin{itemize}
  \item les structures fournisseur ;
  \item les factures fournisseur ;
  \item les clients publics ;
  \item les services destinataires ;
  \item les fichiers d'annuaire normalisés.
\end{itemize}

Le scheduler backend peut exécuter des synchronisations périodiques selon les
variables suivantes :

\begin{lstlisting}
SYNC_CRON_ENABLED=true
SYNC_CRON_INTERVAL_SECONDS=300
SYNC_CRON_MODE=delta
SYNC_CRON_PAGE_SIZE=3000
\end{lstlisting}

\section{Déploiement}

Le projet fournit une stack Docker de production composée de :

\begin{itemize}
  \item un conteneur frontend servi par nginx ;
  \item un conteneur API Python FastAPI ;
  \item un conteneur PostgreSQL ;
  \item un service de migration Alembic ;
  \item des volumes persistants pour la base et l'annuaire.
\end{itemize}

Le démarrage de la stack se fait avec :

\begin{lstlisting}
docker-compose -f docker-compose.prod.yml up -d --build
\end{lstlisting}

Le port HTTP exposé est piloté par \texttt{APP\_HTTP\_PORT}. L'application peut
être publiée sur une IP interne ou un nom DNS interne, par exemple
\texttt{chorusPilot.gigalis.local}.

\section{Tests et validation}

Le projet contient une base de tests orientée à la fois frontend, payloads Chorus,
extraction PDF et automatisation backend.

\begin{itemize}
  \item \textbf{Tests frontend et métier} : exécutés avec Vitest ;
  \item \textbf{Tests Chorus} : validation des payloads, mappers, routes PISTE,
        soumission, extraction et hydratation ;
  \item \textbf{Tests Python} : validation de la configuration, de l'annuaire et
        de la soumission facture ;
  \item \textbf{Guides SOAPUI} : scénarios d'intégration pour tester les appels
        Chorus en environnement de qualification.
\end{itemize}

Commandes utiles :

\begin{lstlisting}
npm run build
npm test
npm run lint
npm run build:test-data
python -m pytest
\end{lstlisting}

\section{Points forts}

\begin{itemize}
  \item Interface unifiée pour consulter, préparer et soumettre les factures ;
  \item séparation claire entre frontend, backend et base de données ;
  \item protection des secrets par passage obligatoire via le backend ;
  \item compatibilité avec les environnements Chorus de qualification et de
        production ;
  \item support de l'extraction PDF via backend Python ;
  \item référentiels clients publics et services destinataires intégrés ;
  \item stack Docker adaptée au déploiement interne ;
  \item jeux de tests et données de démonstration pour sécuriser les évolutions.
\end{itemize}

\section{Axes d'évolution possibles}

Plusieurs améliorations peuvent être envisagées :

\begin{itemize}
  \item gestion multi-utilisateurs complète avec rôles et permissions ;
  \item journal d'audit renforcé par utilisateur réel ;
  \item supervision plus détaillée des jobs de synchronisation ;
  \item notifications en cas de rejet Chorus ou d'échec technique ;
  \item enrichissement de la validation métier avant soumission ;
  \item export de rapports PDF ou CSV pour le suivi administratif ;
  \item intégration d'un module d'archivage documentaire plus complet.
\end{itemize}

\section{Conclusion}

Chorus Pilot fournit une base applicative complète pour piloter le cycle de
vie des factures Chorus Pro. L'application couvre les besoins clés d'un outil
interne : préparation, extraction, contrôle, soumission, synchronisation,
consultation et diagnostic. Son architecture modulaire facilite les évolutions
futures, en particulier autour de la gestion multi-utilisateurs, de la
traçabilité et de l'exploitation en production.

\end{document}
